% Unit 069 terjemahan Bahasa Indonesia dari chapter5.tex baris 740--954.
% Sumber: Methods of Algebra, Volume 2, commit
% 9a5803ff77dd3257484cb177f851a73770a59dd3 (CC BY 4.0).
% stable-unit-id: o014.aljabr2.chapter5.bicomplex-spectral-sequences-applications
% Cakupan: Bagian 5.6 lengkap; berhenti sebelum Bagian 5.7 pada baris 955.

% segment-id: o014.li2.chapter5.bicomplex-spectral-sequences-applications.h001
\section{Barisan spektral bikompleks dan penerapannya}
\label{sec:double-cplx-ss}
\hypertarget{o014.aljabr2.chapter5.bicomplex-spectral-sequences-applications}{ }
% segment-id: o014.li2.chapter5.bicomplex-spectral-sequences-applications.x001
\index{barisan spektral!dari bikompleks}

% segment-id: o014.li2.chapter5.bicomplex-spectral-sequences-applications.p001
Sepanjang bagian ini, kita mengasumsikan bahwa kategori abelian $\mathcal{A}$
mempunyai jumlah langsung terhitung atau produk terhitung yang diperlukan. Misalkan
$X$ bikompleks pada $\mathcal{A}$, ditulis
$X\in\Obj(\cate{C}^2(\mathcal{A}))$. Kompleks total $\tot_{\oplus}X$
mempunyai dua filtrasi menurun,
% segment-id: o014.li2.chapter5.bicomplex-spectral-sequences-applications.q001
\begin{align*}
	\mathrm{F}^p_{\mathrm{I}} (\tot_{\oplus} X)^n & = \bigoplus_{\substack{i+j=n \\ i \geq p}} X^{i, j}, \\
	\mathrm{F}^q_{\mathrm{II}} (\tot_{\oplus} X)^n & = \bigoplus_{\substack{i+j=n \\ j \geq q}} X^{i, j} .
\end{align*}
Dengan memeriksa setiap komponen $(i,j)$ secara tersendiri, terlihat bahwa
% segment-id: o014.li2.chapter5.bicomplex-spectral-sequences-applications.q002
\[ \bigcap_p \mathrm{F}^p_{\mathrm{I}} = 0 = \bigcap_q \mathrm{F}^q_{\mathrm{II}}, \quad \bigcup_p \mathrm{F}^p_{\mathrm{I}} = \tot_{\oplus} X = \bigcup_q \mathrm{F}^q_{\mathrm{II}}. \]

% segment-id: o014.li2.chapter5.bicomplex-spectral-sequences-applications.p002
Dengan demikian diperoleh dua kompleks terfiltrasi. Barisan spektral yang
bersesuaian berturut-turut ditulis
$\mathscr{E}_{\mathrm{I}}=\mathscr{E}_{\mathrm{I}}(X)$ dan
$\mathscr{E}_{\mathrm{II}}=\mathscr{E}_{\mathrm{II}}(X)$, atau secara lebih
konkret sebagai $E_{\mathrm{I},r}^{p,q}$ dan
$E_{\mathrm{II},r}^{p,q}$, dengan $r\in\Z_{\geq0}$. Jika
$X^{p,q}\neq0\implies p,q\geq0$, kita mengatakan bahwa $X$ terletak di
kuadran pertama. Dalam hal ini, untuk
$\star\in\{\mathrm{I},\mathrm{II}\}$ berlaku
% segment-id: o014.li2.chapter5.bicomplex-spectral-sequences-applications.q003
\[ \mathrm{F}_\star^0 \left(\tot_{\oplus} X \right)^n = \left(\tot_{\oplus} X\right)^n, \quad \mathrm{F}_\star^{n+1} \left(\tot_{\oplus} X \right)^n = 0 , \]
sehingga barisan spektral yang bersesuaian juga terletak di kuadran pertama.

% segment-id: o014.li2.chapter5.bicomplex-spectral-sequences-applications.p003
Untuk bikompleks rantai $X$, dengan cara yang sama dapat didefinisikan dua
filtrasi menaik,
% segment-id: o014.li2.chapter5.bicomplex-spectral-sequences-applications.q004
\begin{align*}
	\mathrm{F}_{\mathrm{I}, p} \left( \tot_{\oplus} X \right)_n & = \bigoplus_{\substack{i+j=n \\ i \leq p}} X_{i, j}, \\
	\mathrm{F}_{\mathrm{II}, q} \left( \tot_{\oplus} X \right)_n & = \bigoplus_{\substack{i+j=n \\ j \leq q}} X_{i, j},
\end{align*}
beserta barisan spektral homologis bergradasi ganda yang bersesuaian.

% segment-id: o014.li2.chapter5.bicomplex-spectral-sequences-applications.pr001
\begin{proposition}\label{prop:double-cplx-ss}
	Untuk $X\in\Obj(\cate{C}^2(\mathcal{A}))$, halaman-halaman awal kedua
	barisan spektralnya beserta morfisme-morfisme $d$ yang bersesuaian dapat
	dideskripsikan sebagai berikut:
% segment-id: o014.li2.chapter5.bicomplex-spectral-sequences-applications.tb001
	\[\begin{array}{|c|c|c|c|c|c|} \hline
		& E_0^{p,q} & d_0^{p,q} & E_1^{p,q} & d_1^{p,q} & E_2^{p,q} \\ \hline
		\mathscr{E}_{\mathrm{I}} & X^{p,q} & (-1)^p \dvert^{p,q} & \Hm^q(X^{p, \bullet}, \dvert) & \Hm^q(\dhori^{p, \bullet}) & \Hm_{\mathrm{I}}\Hm_{\mathrm{II}}(X)^{p,q} \\
		\mathscr{E}_{\mathrm{II}} & X^{q,p} & \dhori^{q,p} & \Hm^q(X^{\bullet, p}, \dhori) & (-1)^q \Hm^q(\dvert^{\bullet, p}) & \Hm_{\mathrm{II}} \Hm_{\mathrm{I}}(X)^{q,p} \\ \hline
	\end{array}\]
	Definisi funktor
	$\Hm_{\mathrm{I}},\Hm_{\mathrm{II}}:\cate{C}^2(\mathcal{A})
	\to\cate{C}^2(\mathcal{A})$ terdapat dalam
	\S\ref{sec:double-cplx-coh}, khususnya pada \eqref{eqn:HIHII}.

	Untuk kompleks rantai dan homologinya, cara serupa memberikan tabel berikut.
% segment-id: o014.li2.chapter5.bicomplex-spectral-sequences-applications.tb002
	\[\begin{array}{|c|c|c|c|c|c|} \hline
		& E^0_{p,q} & d^0_{p,q} & E^1_{p,q} & d^1_{p,q} & E^2_{p,q} \\ \hline
		\mathscr{E}_{\mathrm{I}} & X_{p,q} & (-1)^p \dvert_{p,q} & \Hm_q(X_{p, \bullet}, \dvert) & \Hm_q(\dhori_{p, \bullet}) & \Hm_{\mathrm{I}}\Hm_{\mathrm{II}}(X)_{p,q} \\
		\mathscr{E}_{\mathrm{II}} & X_{q,p} & \dhori_{q,p} & \Hm_q(X_{\bullet, p}, \dhori) & (-1)^q \Hm_q(\dvert_{\bullet, p}) & \Hm_{\mathrm{II}} \Hm_{\mathrm{I}}(X)_{q,p} \\ \hline
	\end{array}\]
\end{proposition}
% segment-id: o014.li2.chapter5.bicomplex-spectral-sequences-applications.pf001
\begin{proof}
	Hasil ini diperoleh dengan verifikasi langsung berdasarkan deskripsi dalam
	Proposisi~\ref{prop:filtered-cplx-ss}; tanda pada $\dvert$ berasal dari
	definisi kompleks total. Perinciannya tidak diulangi.
\end{proof}

% segment-id: o014.li2.chapter5.bicomplex-spectral-sequences-applications.p004
Definisi dan hasil yang sama berlaku bagi $\tot_{\Pi}X$, dengan sifat-sifat
yang sepenuhnya serupa.

% segment-id: o014.li2.chapter5.bicomplex-spectral-sequences-applications.p005
Untuk $X\in\Obj(\cate{C}^2_f(\mathcal{A}))$
(Definisi~\ref{def:Cf-Supp}), kompleks totalnya hanya melibatkan jumlah
langsung hingga. Karena itu tidak perlu dibedakan versi $\oplus$ dan $\Pi$;
keduanya ditulis seragam sebagai $\tot X$.

% segment-id: o014.li2.chapter5.bicomplex-spectral-sequences-applications.th001
\begin{theorem}\label{prop:tot-ss}
	Misalkan $X\in\Obj(\cate{C}^2_f(\mathcal{A}))$. Kedua barisan spektral
	$\mathscr{E}_{\mathrm{I}}$ dan $\mathscr{E}_{\mathrm{II}}$ yang
	bersesuaian terbatas dan konvergen kuat, sedangkan kedua filtrasi terinduksi
	pada $\Hm^n(\tot X)$ hingga untuk setiap $n\in\Z$.
\end{theorem}
% segment-id: o014.li2.chapter5.bicomplex-spectral-sequences-applications.pf002
\begin{proof}
	Dari definisi $\cate{C}^2_f(\mathcal{A})$ terlihat bahwa
	$\mathrm{F}_{\mathrm{I}}^\bullet\left(\tot X\right)^n$ dan
	$\mathrm{F}_{\mathrm{II}}^\bullet\left(\tot X\right)^n$ merupakan
	filtrasi hingga untuk setiap $n$. Terapkan
	Teorema~\ref{prop:classical-convergence}.
\end{proof}

% segment-id: o014.li2.chapter5.bicomplex-spectral-sequences-applications.p006
Berikut beberapa penerapan yang representatif. Karena filtrasi dan barisan
spektral yang terlibat selalu hingga, hasil-hasil ini berlaku pada setiap
kategori abelian\footnote{Lebih tepatnya, keberadaan suku-suku limit
$Z_\infty^{p,q}$, $B_\infty^{p,q}$, $E_\infty^{p,q}$ dan syarat tentang
filtrasi menyeluruh dalam Definisi~\ref{def:filtration-properties} tidak
menimbulkan persoalan.}.

% segment-id: o014.li2.chapter5.bicomplex-spectral-sequences-applications.ex001
\begin{example}\label{eg:double-cplx-tot-ss}
	Sekarang kita buktikan kembali Teorema~\ref{prop:double-cplx-tot} dengan
	barisan spektral: jika morfisme
	$f:X\to Y$ dalam $\cate{C}^2_f(\mathcal{A})$ menginduksi
	$\Hm_{\mathrm{II}}\Hm_{\mathrm{I}}(X)\rightiso
	\Hm_{\mathrm{II}}\Hm_{\mathrm{I}}(Y)$ (atau
	$\Hm_{\mathrm{I}}\Hm_{\mathrm{II}}(X)\rightiso
	\Hm_{\mathrm{I}}\Hm_{\mathrm{II}}(Y)$), maka
	$\tot(f):\tot(X)\to\tot(Y)$ merupakan kuasi-isomorfisme.

	Pertama, andaikan
	$\Hm_{\mathrm{II}}\Hm_{\mathrm{I}}(X)\rightiso
	\Hm_{\mathrm{II}}\Hm_{\mathrm{I}}(Y)$. Morfisme $f$ menginduksi
	morfisme barisan spektral
	$\mathscr{E}_{\mathrm{II}}(X)\to\mathscr{E}_{\mathrm{II}}(Y)$.
	Morfisme ini sudah merupakan isomorfisme pada halaman $E_2$, sehingga juga
	memberikan isomorfisme pada halaman limit $E_\infty$
	(Proposisi~\ref{prop:ss-isom-infty}).

	Berdasarkan konvergensi yang dijamin oleh Teorema~\ref{prop:tot-ss}, morfisme
	$\gr^p\left(\Hm^n\tot(X)\right)\to
	\gr^p\left(\Hm^n\tot(Y)\right)$ yang diinduksi oleh $\tot(f)$ merupakan
	isomorfisme untuk semua $p,n\in\Z$. Filtrasi terinduksi pada $\Hm^n$
	diketahui hingga. Karena itu
	$\Hm^n\tot(f):\Hm^n\tot(X)\to\Hm^n\tot(Y)$ juga merupakan isomorfisme
	(Proposisi~\ref{prop:gr-isom}).

	Jika yang ditinjau adalah barisan spektral $\mathscr{E}_{\mathrm{I}}$,
	argumen serupa menunjukkan bahwa $\tot(f)$ merupakan
	kuasi-isomorfisme dari
	$\Hm_{\mathrm{I}}\Hm_{\mathrm{II}}(X)\rightiso
	\Hm_{\mathrm{I}}\Hm_{\mathrm{II}}(Y)$.
\end{example}

% segment-id: o014.li2.chapter5.bicomplex-spectral-sequences-applications.ex002
\begin{example}[Menurunkan bifunktor]\label{eg:bifunctor-ss}
	Misalkan kategori abelian $\mathcal{A}_1$ dan $\mathcal{A}_2$ mempunyai
	cukup banyak objek injektif (atau projektif), sedangkan bifunktor
	$F:\mathcal{A}_1\times\mathcal{A}_2\to\mathcal{B}$ eksak kiri (atau eksak
	kanan) dalam setiap variabel. Kita membahas kasus eksak kiri. Ambil
	$X_i\in\Obj(\mathcal{A}_i)$ dan pilih resolusi injektif
	$0\to X_i\to I_i^0\to\cdots$; definisikan $I_i^n:=0$ jika $n<0$
	($i=1,2$). Data ini membentuk bikompleks di kuadran pertama
% segment-id: o014.li2.chapter5.bicomplex-spectral-sequences-applications.q005
	\[ Y^{p, q} := F\left( I_1^p, I_2^q \right). \]
	Karena itu, barisan spektral kuadran pertama $\mathscr{E}_{\mathrm{I}}$
	yang bersesuaian memenuhi
% segment-id: o014.li2.chapter5.bicomplex-spectral-sequences-applications.q006
	\[ E_1^{p,q} = \Hm^q\left( F(I_1^p, I_2^\bullet) \right) \Rightarrow \Hm^{p+q}(\tot(Y)), \quad p,q \in \Z. \]
	Ruas kanan $\Hm^{p+q}(\tot Y)$ merupakan nilai bifunktor turunan kanan
	$\mathrm{R}^{p+q}F(X_1,X_2)$; lihat
	Definisi~\ref{def:derived-bifunctor}.

	Sekarang andaikan lebih lanjut bahwa $F$ seimbang
	(Definisi~\ref{def:balanced-functor}); dengan demikian,
	$F(I_1^p,\cdot)$ merupakan funktor eksak. Ini menunjukkan bahwa
	$q\neq0\implies E_1^{p,q}=0$, sedangkan
	$E_1^{p,0}=F(I_1^p,X_2)$. Jadi
	$q\neq0\implies E_2^{p,q}=0$, dan
% segment-id: o014.li2.chapter5.bicomplex-spectral-sequences-applications.q007
	\begin{align*}
		E_2^{p,0} & = \Hm^p\left[ \cdots \to F(I_1^p, X_2) \to F(I_1^{p+1}, X_2) \to \cdots \right] \\
		& = (\mathrm{R}_{\mathrm{I}}^p F)(X_1, X_2), \quad \text{dengan notasi Teorema~\ref{prop:balanced-primer}}.
	\end{align*}
	Secara khusus, barisan spektral mengalami degenerasi pada halaman $E_2$,
	sehingga
% segment-id: o014.li2.chapter5.bicomplex-spectral-sequences-applications.q008
	\begin{align*}
		E_2^{p,q} & = E_\infty^{p,q}, \\
		E_\infty^{p,q} & = \gr^p \Hm^{p+q}\left(\tot(Y)\right) = 0, \quad \text{jika}\; q \neq 0, \\
		E_\infty^{p,0} & = \Hm^p\left(\tot(Y)\right) = \mathrm{R}^p F(X_1, X_2) .
	\end{align*}
	Dengan demikian diperoleh
	$\mathrm{R}^nF(X_1,X_2)\simeq
	(\mathrm{R}_{\mathrm{I}}^nF)(X_1,X_2)$. Jika sebagai gantinya digunakan
	$\mathscr{E}_{\mathrm{II}}$, argumen yang sama memberikan
	$\mathrm{R}^nF(X_1,X_2)\simeq
	(\mathrm{R}_{\mathrm{II}}^nF)(X_1,X_2)$. Jadi
	$\mathrm{R}_{\mathrm{I}}F\simeq\mathrm{R}_{\mathrm{II}}F$; inilah
	pernyataan Teorema~\ref{prop:balanced-primer}.
\end{example}

% segment-id: o014.li2.chapter5.bicomplex-spectral-sequences-applications.ex003
\begin{example}[Barisan spektral funktor hiperturunan]
\label{eg:hyperderived-ss}
% segment-id: o014.li2.chapter5.bicomplex-spectral-sequences-applications.x002
	\index{funktor hiperturunan}
% segment-id: o014.li2.chapter5.bicomplex-spectral-sequences-applications.x003
	\index{barisan spektral!funktor hiperturunan}
	Misalkan $\mathcal{A}$ dan $\mathcal{B}$ kategori abelian,
	$\mathcal{A}$ mempunyai cukup banyak objek injektif (atau projektif), dan
	$F:\mathcal{A}\to\mathcal{B}$ funktor aditif eksak kiri (atau eksak kanan).
	Pada paruh pertama \S\ref{sec:derived-primer}, untuk kompleks
	$X\in\Obj(\cate{C}^+(\mathcal{A}))$ (atau
	$X\in\Obj(\cate{C}^-(\mathcal{A}))$), kita mendefinisikan funktor turunan
	kanan $\mathrm{R}^nF(X)$ (atau funktor turunan kiri $\mathrm{L}_nF(X)$),
	dengan $n\in\Z$. Jika $X\in\Obj(\mathcal{A})$, funktor-funktor tersebut
	merupakan funktor turunan dalam pengertian klasik. Sebaliknya, kasus kompleks
	umum lazim disebut funktor hiperturunan. Kedua kasus ini dihubungkan oleh
	barisan spektral. Kita menjabarkan kasus funktor turunan kanan; pertukaran
	superskrip dengan subskrip memberikan versi funktor turunan kiri.

	Ambil $X\in\Obj(\cate{C}^+(\mathcal{A}))$. Pandang $X$ sebagai bikompleks
	yang terkonsentrasi pada baris ke-$0$, lalu ambil resolusi
	Cartan--Eilenberg $\epsilon:X\to I$ yang diberikan oleh
	Teorema~\ref{prop:CE-resolution}; ini merupakan morfisme dalam
	$\cate{C}^2_f(\mathcal{A})$. Menurut Catatan~\ref{rem:double-cplx-resolution}
	atau hasil Contoh~\ref{eg:double-cplx-tot-ss},
	$\tot(\epsilon):X=\tot(X)\to\tot(I)$ merupakan kuasi-isomorfisme dalam
	$\cate{C}^+(\mathcal{A})$, sehingga merupakan resolusi injektif dari $X$.

	Dari bikompleks
	$(\cate{C}^2F)(I)\in
	\Obj\left(\cate{C}^2_f(\mathcal{B})\right)$, bangun barisan spektral
	konvergen $\mathscr{E}_{\mathrm{I}}$ dan $\mathscr{E}_{\mathrm{II}}$.
	Keduanya konvergen menuju sasaran yang sama,
	$\Hm^{p+q}\tot\left((\cate{C}^2F)I\right)
	=\Hm^{p+q}\cate{C}F(\tot I)$, yaitu nilai funktor hiperturunan
	$\mathrm{R}^{p+q}F(X)$.

	Pertama tinjau $\mathscr{E}_{\mathrm{I}}$. Karena $I^{p,\bullet}$
	merupakan resolusi injektif dari $X^p$, berlaku
	$E_{\mathrm{I},1}^{p,q}=\Hm^q\left(FI^{p,\bullet}\right)
	\simeq\mathrm{R}^qF(X^p)$; ruas kanan merupakan nilai funktor turunan klasik
	$\mathrm{R}^qF$ pada $X^p$.\footnote{Koreksi penerjemah: sumber menuliskan
	$\mathrm{R}^pF$, sedangkan rumus sebelumnya dan indeks kohomologi vertikal
	menuntut $\mathrm{R}^qF$.} Dengan alasan serupa,
% segment-id: o014.li2.chapter5.bicomplex-spectral-sequences-applications.q009
	\[ E_{\mathrm{I}, 2}^{p,q} = \Hm^p\left[ \cdots \to \mathrm{R}^q F(X^p) \to \mathrm{R}^q F(X^{p+1}) \to \cdots \right]. \]

	Yang mungkin lebih menarik adalah halaman $E_2$ dari
	$\mathscr{E}_{\mathrm{II}}$. Berdasarkan sifat resolusi Cartan--Eilenberg,
	hasil pengambilan kohomologi horizontal
	$\Hm_{\mathrm{I}}(I)^{q,\bullet}$ memberikan resolusi injektif dari
	$\Hm^q(X)$, sedangkan Catatan~\ref{rem:CE-split} menunjukkan bahwa $F$
	mempertahankan kohomologi horizontal:
	$\Hm_{\mathrm{I}}(\cate{C}^2F(I))^{q,\bullet}
	\simeq\cate{C}F\left(\Hm_{\mathrm{I}}(I)^{q,\bullet}\right)$. Karena itu,
% segment-id: o014.li2.chapter5.bicomplex-spectral-sequences-applications.g001
	\begin{align*}
		E_{\mathrm{II},2}^{p, q} & =
		\begin{tikzpicture}[scale=0.7, baseline=(O)]
			\draw[-Latex] (0, -1) -- (0, 1) node[right] {$p$} node[left=0.5em] {\scriptsize\text{ambil $\Hm$ kemudian}};
			\draw[-Latex] (-1, 0) -- (1, 0) node[below] {$q$} node[right=0.5em] {\scriptsize\text{ambil $\Hm$ dahulu}};
			\coordinate (O) at (0, -0.3);
			\draw (current bounding box.north west) rectangle ([yshift=-1ex] current bounding box.south east);
		\end{tikzpicture}
		= (\mathrm{R}^p F)\left( \Hm^q(X)\right) \\
		& \Rightarrow \mathrm{R}^{p+q}F(X) , \quad p, q \in \Z.
	\end{align*}
\end{example}

% segment-id: o014.li2.chapter5.bicomplex-spectral-sequences-applications.p007
Barisan spektral Grothendieck yang diperkenalkan berikut berkenaan dengan
funktor turunan dari suatu komposisi funktor. Barisan ini mencakup suatu kelas
besar barisan spektral dalam
geometri dan aljabar, dan argumennya, seperti dalam
Contoh~\ref{eg:hyperderived-ss}, juga didasarkan pada resolusi
Cartan--Eilenberg.

% segment-id: o014.li2.chapter5.bicomplex-spectral-sequences-applications.th002
\begin{theorem}[Barisan spektral Grothendieck]
\label{prop:Grothendieck-ss}
% segment-id: o014.li2.chapter5.bicomplex-spectral-sequences-applications.x004
	\index{barisan spektral!Grothendieck}
	Tinjau funktor-funktor aditif antarkategori abelian
% segment-id: o014.li2.chapter5.bicomplex-spectral-sequences-applications.q010
	\[ \mathcal{A} \xrightarrow{F} \mathcal{A}' \xrightarrow{F'} \mathcal{A}''. \]
	Misalkan keduanya funktor eksak kiri (atau eksak kanan), $\mathcal{A}$ dan
	$\mathcal{A}'$ mempunyai cukup banyak objek injektif (atau projektif), dan
	$F$ memetakan objek injektif (atau projektif) ke objek $F'$-asiklik dalam
	pengertian Konvensi~\ref{con:F-acyclic}. Maka, untuk setiap
	$X\in\Obj(\mathcal{A})$, terdapat barisan spektral kuadran pertama
	kohomologis (atau homologis) bergradasi ganda
% segment-id: o014.li2.chapter5.bicomplex-spectral-sequences-applications.q011
	\begin{align*}
		E_2^{p,q} = (\mathrm{R}^p F') (\mathrm{R}^q F)(X) & \Rightarrow \mathrm{R}^{p+q} (F'F)(X), \\
		\text{atau} \quad E^2_{p,q} = (\mathrm{L}_p F') (\mathrm{L}_q F)(X) & \Rightarrow \mathrm{L}_{p+q}(F'F)(X).
	\end{align*}
	Sifat konvergensinya adalah seperti dalam Teorema~\ref{prop:tot-ss}.
	Barisan eksak suku berderajat rendah yang bersesuaian
	(Korolari~\ref{prop:low-degree-ss}) masing-masing dapat ditulis sebagai
% segment-id: o014.li2.chapter5.bicomplex-spectral-sequences-applications.q012
	\begin{equation*}\begin{split}
		0 \to (\mathrm{R}^1 F')(FX) & \to \mathrm{R}^1(F'F)(X) \\
		& \to F'\left( (\mathrm{R}^1 F)X \right) \to (\mathrm{R}^2 F')(FX) \to \mathrm{R}^2(F'F)(X),
	\end{split}\end{equation*}
	atau
% segment-id: o014.li2.chapter5.bicomplex-spectral-sequences-applications.q013
	\begin{equation*}\begin{split}
		\mathrm{L}_2(F'F)(X) \to (\mathrm{L}_2 F')(FX) & \to F'\left( (\mathrm{L}_1 F)X \right) \\
		& \to \mathrm{L}_1(F'F)(X) \to (\mathrm{L}_1 F')(FX) \to 0.
	\end{split}\end{equation*}
\end{theorem}
% segment-id: o014.li2.chapter5.bicomplex-spectral-sequences-applications.pf003
\begin{proof}
	Berdasarkan dualitas, cukup dibahas kasus kohomologis. Ambil resolusi
	injektif $0\to X\to I^0\to I^1\to\cdots$ dari $X$ dan definisikan
	$I^n:=0$ untuk $n<0$. Dalam $\mathcal{A}'$, ambil resolusi
	Cartan--Eilenberg $\cate{C}F(I)\to J$ dari
	$\cate{C}F(I):=(FI^p)_p$ (Teorema~\ref{prop:CE-resolution}), lalu tinjau
	bikompleks kuadran pertama
	$\cate{C}^2F'(J)=(F'(J^{p,q}))_{p,q}$ beserta barisan spektral konvergen
	$\mathscr{E}_{\mathrm{I}}$ dan $\mathscr{E}_{\mathrm{II}}$ yang
	bersesuaian. Pertama,
% segment-id: o014.li2.chapter5.bicomplex-spectral-sequences-applications.g002
	\begin{align*}
		E_{\mathrm{I},2}^{p, q} & =
		\begin{tikzpicture}[scale=0.7, baseline=(O)]
			\draw[-Latex] (0, -1) -- (0, 1) node[right] {$q$} node[left=0.5em] {\scriptsize\text{ambil $\Hm$ dahulu}};
			\draw[-Latex] (-1, 0) -- (1, 0) node[below] {$p$} node[right=0.5em] {\scriptsize\text{ambil $\Hm$ kemudian}};
			\coordinate (O) at (0, -0.3);
			\draw (current bounding box.north west) rectangle ([yshift=-1ex] current bounding box.south east);
		\end{tikzpicture} \\
		& = \Hm^p\left[ \cdots \to (\mathrm{R}^q F')(FI^p) \to (\mathrm{R}^q F')(FI^{p+1}) \to \cdots \right] \\
		& \Rightarrow \Hm^{p+q}\tot\left(\cate{C}^2 F'(J)\right).
	\end{align*}

	Karena menurut hipotesis $F(I^p)$ merupakan objek $F'$-asiklik, berlaku
	$E_{\mathrm{I},2}^{p,q}=0$ jika $q\neq0$. Argumen yang persis sama dengan
	Contoh~\ref{eg:bifunctor-ss} kemudian menunjukkan bahwa
	$\mathscr{E}_{\mathrm{I}}$ mengalami degenerasi pada halaman $E_2$, dan
% segment-id: o014.li2.chapter5.bicomplex-spectral-sequences-applications.q014
	\[ E_{\mathrm{I}, \infty}^{p,q} = E_{\mathrm{I}, 2}^{p,q} = \begin{cases}
		\Hm^p\left( F'F(I^\bullet) \right) = \mathrm{R}^p (F' F)(X), & q = 0, \\
		0, & q \neq 0,
	\end{cases}\]
	serta
	$\Hm^p\tot\left(\cate{C}^2F'(J)\right)
	=E_{\mathrm{I},\infty}^{p,0}=\mathrm{R}^p(F'F)(X)$.

	Sekarang tinjau $\mathscr{E}_{\mathrm{II}}$. Tekniknya serupa dengan
	Contoh~\ref{eg:hyperderived-ss}: kolom ke-$q$ dari kohomologi horizontal,
	$\Hm_{\mathrm{I}}(J)^{q,\bullet}$, memberikan resolusi injektif dari
	$\Hm^q(\cate{C}F(I))=(\mathrm{R}^qF)(X)$. Dengan mengingat pula bahwa $F'$
	mempertahankan kohomologi horizontal, diperoleh
% segment-id: o014.li2.chapter5.bicomplex-spectral-sequences-applications.g003
	\begin{align*}
		E_{\mathrm{II},2}^{p, q} & =
		\begin{tikzpicture}[scale=0.7, baseline=(O)]
			\draw[-Latex] (0, -1) -- (0, 1) node[right] {$p$} node[left=0.5em] {\scriptsize\text{ambil $\Hm$ kemudian}};
			\draw[-Latex] (-1, 0) -- (1, 0) node[below] {$q$} node[right=0.5em] {\scriptsize\text{ambil $\Hm$ dahulu}};
			\coordinate (O) at (0, -0.3);
			\draw (current bounding box.north west) rectangle ([yshift=-1ex] current bounding box.south east);
		\end{tikzpicture}
		= (\mathrm{R}^p F')(\mathrm{R}^q F)\left( X \right) \\
		& \Rightarrow \Hm^{p+q}\tot\left(\cate{C}^2 F'(J)\right) = \mathrm{R}^{p+q}(F'F)(X).
	\end{align*}

	Jadi $\mathscr{E}_{\mathrm{II}}$ merupakan barisan spektral yang dikehendaki
	untuk bagian pertama. Dengan memasukkan deskripsi
	$E_{\mathrm{II},2}^{p,q}$ ke barisan eksak suku berderajat rendah dari
	$\mathscr{E}_{\mathrm{II}}$ (Korolari~\ref{prop:low-degree-ss}), diperoleh
	langsung pernyataan bagian terakhir.
\end{proof}

% segment-id: o014.li2.chapter5.bicomplex-spectral-sequences-applications.p008
Barisan spektral Grothendieck dapat dipandang sebagai versi konkret dari
Teorema~\ref{prop:derived-composite}, tetapi informasi yang dimuatnya lebih
rinci daripada isomorfisme dalam kategori turunan. Kita akan memberikan contoh
penerapannya dalam \S\ref{sec:LHS-SS}.

% segment-id: o014.li2.chapter5.bicomplex-spectral-sequences-applications.ex004
\begin{example}[Perubahan gelanggang]\label{eg:change-of-rings-SS}
% segment-id: o014.li2.chapter5.bicomplex-spectral-sequences-applications.x005
	\index{perubahan gelanggang}
% segment-id: o014.li2.chapter5.bicomplex-spectral-sequences-applications.x006
	\index{barisan spektral!perubahan gelanggang}
	Misalkan $R\to S$ homomorfisme gelanggang. Ambil $X$ sebagai modul kanan
	$S$ dan $Y$ sebagai modul kiri $R$. Tuliskan $S_R$ (atau $X_R$) untuk $S$
	(atau $X$) yang dipandang sebagai modul kanan $R$. Akan ditunjukkan bahwa
	terdapat barisan spektral kuadran pertama homologis bergradasi ganda
% segment-id: o014.li2.chapter5.bicomplex-spectral-sequences-applications.q015
	\begin{gather*}
		E^2_{p,q} = \Tor^S_p\left( X, \Tor^R_q(S_R, Y) \right) \Rightarrow \Tor^R_{p+q}(X_R, Y);
	\end{gather*}
	Di sini $\Tor^R_q(S_R,Y)$ diberi struktur modul kiri $S$ seperti dalam
	Catatan~\ref{rem:ExtTor-bimodule}. Ini merupakan penerapan langsung
	Teorema~\ref{prop:Grothendieck-ss}, yang berasal dari isomorfisme funktor
% segment-id: o014.li2.chapter5.bicomplex-spectral-sequences-applications.q016
	\[ X \dotimes{S} \left( S \dotimes{R} (\cdot) \right) \simeq X \dotimes{R} (\cdot) : R\dcate{Mod} \to \cate{Ab}, \]
	dengan funktor dalam
	$S\dotimes{R}(\cdot):R\dcate{Mod}\to S\dcate{Mod}$ mempunyai funktor
	turunan kiri $\Tor^R_\bullet(S_R,\cdot)$ yang bernilai dalam
	$S\dcate{Mod}$. Berdasarkan isomorfisme di atas, funktor dalam memetakan
	modul datar ke modul datar, sehingga versi homologis barisan spektral
	Grothendieck memang dapat diterapkan.

	Dengan notasi dan teknik yang sepenuhnya serupa, untuk modul kanan $R$, $X$,
	dan modul kiri $S$, $Y$, juga diperoleh
% segment-id: o014.li2.chapter5.bicomplex-spectral-sequences-applications.q017
	\[ E^2_{p,q} = \Tor^S_p\left( \Tor^R_q(X, {}_R S), Y \right) \Rightarrow \Tor^R_{p+q}(X, {}_R Y). \]

	Selanjutnya tinjau funktor $\Ext$. Ambil $X$ sebagai modul kiri $S$ dan
	$Y$ sebagai modul kiri $R$, serta tuliskan ${}_RX$ untuk $X$ yang dipandang
	sebagai modul kiri $R$. Berikan $\Ext_R^q({}_RS,Y)$ struktur modul kiri $S$
	seperti dalam Catatan~\ref{rem:ExtTor-bimodule}. Maka terdapat barisan
	spektral kuadran pertama kohomologis bergradasi ganda
% segment-id: o014.li2.chapter5.bicomplex-spectral-sequences-applications.q018
	\begin{align*}
		E_2^{p,q} = \Ext_S^p\left( X, \Ext_R^q({}_R S, Y) \right) & \Rightarrow \Ext_R^{p+q}({}_R X, Y), \\
		E_2^{p,q} = \Ext_S^p\left( \Tor^R_q(S_R, Y), X \right) & \Rightarrow \Ext_R^{p+q}(Y, {}_R X).
	\end{align*}
	Kedua barisan itu berturut-turut bersesuaian dengan isomorfisme komposisi funktor
% segment-id: o014.li2.chapter5.bicomplex-spectral-sequences-applications.q019
	\begin{gather*}
		\Hom_S\left( X, \Hom_R({}_R S, \cdot) \right) \simeq \Hom_R({}_R X, \cdot), \\
		\Hom_S\left( S \dotimes{R} (\cdot), X \right) \simeq \Hom_R(\cdot, {}_R X),
	\end{gather*}
	yaitu adjungsi-adjungsi yang diperkenalkan dalam
	\cite[Korolari 6.6.8]{Li1}. Pemakaian campuran funktor turunan kanan
	$\Ext_S^p$ dan funktor turunan kiri $\Tor^R_q$ dalam baris kedua tidak
	menimbulkan masalah, sebab variabel pertama $\Hom_S$ berada dalam kategori
	berlawanan $S\dcate{Mod}^{\opp}$.
\end{example}

% segment-id: o014.li2.chapter5.bicomplex-spectral-sequences-applications.p009
Barisan-barisan spektral di atas sebaiknya dibandingkan dengan versi kategori
turunan dalam \S\ref{sec:otimesL}.

% segment-id: o014.li2.chapter5.bicomplex-spectral-sequences-applications.ex005
\begin{example}\label{eg:change-of-rings-SS-bis}
% segment-id: o014.li2.chapter5.bicomplex-spectral-sequences-applications.x007
\index{teorema koefisien universal}
	Untuk barisan spektral
	$E_2^{p,q}=\Ext_S^p\left(\Tor^R_q(S_R,Y),X\right)
	\Rightarrow\Ext_R^{p+q}(Y,{}_RX)$ dari
	Contoh~\ref{eg:change-of-rings-SS}, kita catat sebuah kasus khusus yang
	berguna, sekaligus sebagai latihan. Andaikan dimensi global
	$S\dcate{Mod}$ tidak lebih dari $1$
	(Definisi--Proposisi~\ref{def:global-dim}). Akibatnya, suku-suku taknol
	$E_2^{p,q}$ terkonsentrasi pada wilayah $p=0,1$. Jadi, untuk $n\geq0$,
	filtrasi menurun $\mathrm{F}^\bullet$ pada $\Ext_R^n(Y,{}_RX)$ harus
	berbentuk
% segment-id: o014.li2.chapter5.bicomplex-spectral-sequences-applications.q020
	\[ \Ext_R^n(Y, {}_R X) = \mathrm{F}^0 \underbracket{\supset}_{E_\infty^{0, n}} \mathrm{F}^1 \underbracket{\supset}_{E_\infty^{1, n-1}} \mathrm{F}^2 = \{0\}. \]

	Selain itu, Contoh~\ref{eg:SS-finite-width} menunjukkan bahwa barisan
	spektral mengalami degenerasi pada halaman $E_2$, sehingga
	$E_2^{p,q}\simeq E_\infty^{p,q}$. Dengan menggabungkan semuanya diperoleh
	barisan eksak pendek
% segment-id: o014.li2.chapter5.bicomplex-spectral-sequences-applications.q021
	\[ 0 \to \Ext_S^1\left( \Tor^R_{n-1}(S_R, Y), X \right) \to \Ext_R^n(Y, {}_R X) \to \Hom_S\left( \Tor^R_n(S_R, Y), X \right) \to 0; \]
	dengan konvensi $\Tor^R_{-1}=0$. Barisan di atas dapat dipandang sebagai
	salah satu varian teorema koefisien universal.

	Dengan cara serupa, andaikan modul kanan $S$, $X$, mempunyai dimensi $\Tor$
	tidak lebih dari $1$ (Contoh~\ref{eg:Tor-dimension}). Barisan spektral
	$E^2_{p,q}=\Tor^S_p(X,\Tor^R_q(S_R,Y))
	\Rightarrow\Tor^R_{p+q}(X_R,Y)$ mengalami degenerasi pada halaman $E^2$.
	Dengan memeriksa filtrasi menaik yang bersesuaian pada
	$\Tor^R_n(X_R,Y)$, diperoleh barisan eksak pendek
% segment-id: o014.li2.chapter5.bicomplex-spectral-sequences-applications.q022
	\[ 0 \to X \dotimes{S} \Tor^R_n(S_R, Y) \to \Tor^R_n(X_R, Y) \to \Tor^S_1(X, \Tor^R_{n-1}(S_R, Y)) \to 0. \]

	Jika $S$ merupakan domain ideal utama, baik hipotesis dimensi global maupun
	dimensi $\Tor$ otomatis terpenuhi.
\end{example}
